This section describes the first release (v2.0.0), which ran the research
MobileNetV3-Large; Section~\ref{sec:phone} covers what changed after it met real
photos. The app runs the model entirely on the phone, so photos never leave the device
and no network connection is needed; the release build does not even request
internet access. We export the deployed MobileNetV3-Large to ONNX in 32-bit
floating point (\AppModelMb{}\,MB) and run it with ONNX Runtime inside a Flutter
app. Dynamic 8-bit quantisation would cut the file to 5.1\,MB, but it changed
many of the model's predictions, so we did not use it. The model ships with a
metadata file holding the label order, preprocessing constants and OOD
threshold, which keeps the app and the training code from drifting apart. The
app labels the grade and shelf-life fields as estimates and avoids food-safety
wording (``looks fresh'' rather than ``safe to eat'').

We checked the port on \AppParityN{} held-out test images run through the full
on-device pipeline, from JPEG or PNG decoding to the gate decision. Freshness,
produce type and the gate decision agreed with PyTorch on every image; the
largest logit difference, \AppParityMaxDiff{}, comes from the JPEG decoder rather
than the model. The arm64 APK is \AppApkArm{}\,MB (\AppApkArmSeven{}\,MB for
32-bit phones), of which the ONNX Runtime library and the model take
\AppOrtMb{} and \AppModelMb{}\,MB. On an Android emulator with two virtual
cores, preprocessing and inference take \AppPipeMs{}\,ms (median), a full scan
in the app about \AppScanMs{}\,ms, and a cold start \AppColdMs{}\,s. Memory use
is \AppRamIdle{}\,MB at idle and \AppRamLoaded{}\,MB once the model is loaded,
peaking at \AppRamPeak{}\,MB during scans, and each saved scan adds about
\AppScanKb{}\,kB of storage. Emulator timings are only a guide to phone
performance, which we did not measure.
